Read in Gauge Theory Gravity, the metric below is the derived metric, $\nabla$ is the covariant derivative $D$ on physical vectors, and the Riemann and Weyl tensors are read from the field strength and its Weyl map (Fact \claimref{synthesis.imported-gtg-dictionary}, clauses (i), (ii), (iv) and (v)). \emph{Conventions.} Signature $(-,+,+,+)$; $R^{a}{}_{bcd} = \partial_{c}\Gamma^{a}_{db} - \partial_{d}\Gamma^{a}_{cb} + \Gamma^{a}_{ce}\Gamma^{e}_{db} - \Gamma^{a}_{de}\Gamma^{e}_{cb}$, so that $R(e_{1},e_{2},e_{1},e_{2})$ is the sectional curvature of an orthonormal plane and $\Ric^{\circ}_{bd} = R^{a}{}_{bad}$ is positive on round spheres; this $\Ric^{\circ}$ is the one in the Raychaudhuri equation of Fact \claimref{synthesis.imported-raychaudhuri-null-congruence}. The element $\Sigma$, its null normals $k$, $l$ with $g(k,l) = -1$, the generators of $\mathcal{H}$ with affine parameter $\lambda$, the cuts $\Sigma_{\lambda}$, the frame $e_{a}$ (Lie-dragged along $k$, so $\nabla_{k}e_{a} = \nabla_{e_{a}}k$), $B^{(k)}$, $B^{(l)}$, $\omega$, $\widehat{C}_{k}$, $Z$ and $P$ are as in the statement; screen indices are raised with $h$. The generators of a null hypersurface are twist-free.

\emph{Step 1 (Gauss equation).} The normal part of $\nabla_{e_{a}}e_{b}$ is $K_{ab} = B^{(l)}_{ab}\,k + B^{(k)}_{ab}\,l$: indeed $g(K_{ab}, k) = g(\nabla_{e_{a}}e_{b}, k) = -g(e_{b}, \nabla_{e_{a}}k) = -B^{(k)}_{ab}$, which is $g(B^{(k)}_{ab}\,l, k)$ since $g(l,k) = -1$, and likewise with $k$ and $l$ exchanged; both forms are symmetric because $[e_{a}, e_{b}]$ is tangent to $\Sigma$. The Gauss equation ${}^{(2)}\!R_{abcd} = R_{abcd} + g(K_{ac},K_{bd}) - g(K_{ad},K_{bc})$ (standard, cited from memory) with $g(K_{ac},K_{bd}) = -(B^{(l)}_{ac}B^{(k)}_{bd} + B^{(k)}_{ac}B^{(l)}_{bd})$, contracted with $h^{ac}h^{bd}$, gives
\[
{}^{(2)}\!R = P - 2\,\theta_{k}\theta_{l} + 2\,B^{(k)}_{ab}B^{(l)ab} .
\]

\emph{Step 2 (the projected part).} Along the generator through $p$, $\partial_{\lambda}h_{ab} = 2B^{(k)}_{(ab)} = 0$ at $p$ and $\partial_{\lambda}\sqrt{h} = \theta_{k}\sqrt{h} = 0$ at $p$, so only the curvature and the frame move at first order. Since $g(\nabla_{e_{a}}k, k) = 0$ and $g(\nabla_{e_{a}}k, l) = -\omega_{a}$, $\nabla_{e_{a}}k = B^{(k)}{}_{a}{}^{b}e_{b} + \omega_{a}k$, which is $\omega_{a}k$ at $p$. Hence at $p$
\[
\frac{dP}{d\lambda} = (\nabla_{k}R)(e_{a},e_{b},e^{a},e^{b}) + h^{ac}h^{bd}\big(\omega_{a}R(k,e_{b},e_{c},e_{d}) + \omega_{b}R(e_{a},k,e_{c},e_{d}) + \omega_{c}R(e_{a},e_{b},k,e_{d}) + \omega_{d}R(e_{a},e_{b},e_{c},k)\big),
\]
and the pair symmetry and antisymmetries of $R$ make the four frame terms equal, each $\omega^{a}h^{bd}R(k,e_{b},e_{a},e_{d}) = \omega^{a}Z_{a}$ (since $R(k,e_{b},e_{a},e_{b}) = R(e_{b},k,e_{b},e_{a})$). This is the second formula of (i). The second Bianchi identity $\nabla_{k}R_{abab} + \nabla_{a}R_{abbk} + \nabla_{b}R_{abka} = 0$, summed over the screen indices, gives $(\nabla_{k}R)(e_{a},e_{b},e^{a},e^{b}) = 2(\nabla_{e^{b}}R)(e_{a},k,e^{a},e_{b})$, the divergence form of the gradient term.

\emph{Step 3 (the extrinsic part).} At $p$, $B^{(k)} = 0$, so the derivative of $-2\theta_{k}\theta_{l} + 2B^{(k)}_{ab}B^{(l)ab}$ needs only $dB^{(k)}/d\lambda$ at $p$, not $dB^{(l)}/d\lambda$. By Fact \claimref{synthesis.imported-raychaudhuri-null-congruence}, at $p$ (where $\theta_{k} = \sigma_{k} = 0$) $d\theta_{k}/d\lambda = -\Ric^{\circ}(k,k)$ and $d\sigma_{k}/d\lambda = -\widehat{C}_{k}$, so $dB^{(k)}/d\lambda = -\widehat{C}_{k} - \tfrac12\Ric^{\circ}(k,k)\,h$. Therefore
\[
\frac{d}{d\lambda}\big(-2\theta_{k}\theta_{l} + 2B^{(k)}_{ab}B^{(l)ab}\big) = 2\Ric^{\circ}(k,k)\theta_{l} - 2\widehat{C}_{k\,ab}B^{(l)ab} - \Ric^{\circ}(k,k)\theta_{l} = \theta_{l}\Ric^{\circ}(k,k) - 2\sigma_{l}^{ab}\widehat{C}_{k\,ab},
\]
using that $\widehat{C}_{k}$ is trace-free on the screen. With Steps 1 and 2 and $\partial_{\lambda}\sqrt{h} = 0$ at $p$ this is the first formula of (i). Only $B^{(k)}$, $B^{(l)}$ and $\omega$ at $p$ and the curvature and its first derivatives at $p$ enter; third- and higher-order data of the embedding of $\Sigma$, and the normalisation of $k$ beyond its first derivative, do not.

\emph{Step 4 (free data).} In coordinates in which $p = 0$, the metric is $-2\,du\,dv + dx^{2} + dy^{2}$ at $p$ with vanishing first derivatives, $k = \partial_{u}$ and $l = \partial_{v}$ at $p$, an element $X(s_{1},s_{2}) = (u(s), v(s), s_{1} + \dots, s_{2} + \dots)$ has $B^{(k)}_{ab}(p) = \partial_{a}\partial_{b}v(0)$ and $B^{(l)}_{ab}(p) = \partial_{a}\partial_{b}u(0)$. So $v$ without quadratic terms makes the element admissible, and the quadratic part of $u$ prescribes any symmetric $B^{(l)}(p)$. Rescaling $k \mapsto e^{f}k$ by a function on $\Sigma$ constant along the generators, with $f(p) = 0$, keeps the generators affinely parametrised, keeps $B^{(k)}(p) = 0$ and $B^{(l)}(p)$ (with $l \mapsto e^{-f}l$), and changes $\omega$ to $\omega + df$, so every covector $\omega(p)$ is realised. This proves the last sentence of (i).

\emph{Step 5 (Clausius at first order).} In the setting of Fact \claimref{synthesis.imported-jacobson-clausius-1995}, per unit cross-section area in the thin-pencil limit (H2), with the patch $\lambda \in [\lambda_{0}, 0]$ and $\chi = -\kappa\lambda k$ up to $O(\lambda^{2})$, Theorem \claimref{synthesis.einstein-equations-from-c-loc-and-clausius}(2) gives the heat $\delta Q = \tfrac12\kappa T(k,k)\lambda_{0}^{2} + O(\lambda_{0}^{3})$ and the area change $\delta A = \tfrac12\Ric^{\circ}(k,k)\lambda_{0}^{2} + O(\lambda_{0}^{3})$, both without a term of order $\lambda_{0}$. The Euler part of the Jacobson-Myers entropy (Fact \claimref{synthesis.imported-wald-noether-charge-entropy}(c)), $\frac{\alpha}{2G}\int{}^{(2)}\!R\,dA$, changes per unit area by $\frac{\alpha}{2G}\big[\sqrt{h}\,{}^{(2)}\!R\big]_{\lambda_{0}}^{0}/\sqrt{h}(0) = -\frac{\alpha}{2G}\,r\,\lambda_{0} + O(\lambda_{0}^{2})$, with $r$ the first rate of (i); likewise with $P$ and the second rate. With $T_{\chi} = \kappa/2\pi$,
\[
\delta Q - T_{\chi}\,\delta S = \frac{\kappa}{2\pi}\,\frac{\alpha}{2G}\,r\,\lambda_{0} + O(\lambda_{0}^{2}) .
\]
(H1) holds at leading order on every admissible patch, so $\alpha r = 0$, and $\alpha = -3/(4\Lambda) \neq 0$ forces $r = 0$ for every admissible element through $p$.

\emph{Step 6 (what $r = 0$ for all elements means).} For the Jacobson-Myers form, $r$ is affine in the free data $(\theta_{l}, \sigma_{l}, \omega)$ of Step 4, with coefficients $\Ric^{\circ}(k,k)$, $-2\widehat{C}_{k}$ and $4Z$ and constant term the gradient term. It vanishes for all data if and only if all four vanish. $\Ric^{\circ}(k,k) = 0$ for every null $k$ is the vanishing of the trace-free Ricci tensor (a symmetric form vanishing on the null cone is a multiple of $g$; the check confirms rank nine on eleven rational null vectors). $\widehat{C}_{k} = 0$ for every null $k$ and its screen is the vanishing of the Weyl tensor (the check builds the ten-dimensional space of Weyl tensors as the trace-free part of the span of Kulkarni-Nomizu products and finds the map to these components injective, rank ten). Then the curvature at $p$ is $K(g_{ac}g_{bd} - g_{ad}g_{bc})$, for which $Z_{b} = K(h^{ac}g_{ac}g(k,e_{b}) - h^{ac}g(e_{a},e_{b})g(k,e_{c})) = 0$. If this holds at every point of a connected region, the contracted Bianchi identity makes $K$ constant (Schur's lemma, standard, cited from memory), so $\nabla R = 0$ and the gradient term vanishes as well; conversely in constant curvature all four coefficients vanish. For the projected form, $r = 4\omega\cdot Z + (\nabla_{k}R)(e,e,e,e)$, and $r = 0$ for all $\omega$ requires $Z = 0$ and the gradient term to vanish, for every null $k$ and screen.

\emph{Step 7 (examples).} In the plane waves, the gravitational fields $\hb^{-1}(a) = a + \tfrac12 H\,(a\cdot e_{v})\,e_{v}$ ($e_{v}$ null, $e_{u}\cdot e_{v} = -1$) with derived metric $-2\,du\,dv + H\,du^{2} + dx^{2} + dy^{2}$, the only curvature components are $R_{uiuj} = -\tfrac12\partial_{i}\partial_{j}H$. For $k = \partial_{u}$, $Z = 0$, and for $H$ quadratic the gradient vanishes, so $r = \theta_{l}\Ric^{\circ}(k,k) - 2\sigma_{l}\!:\!\widehat{C}_{k}$: for the vacuum wave $H = 3x^{2} - 3y^{2} + 2xy$ an element with $\sigma_{l} \neq 0$ gives $r \neq 0$ ($32$ in the check), so the Jacobson-Myers form fails (H1) in vacuum; for the null dust $H = x^{2} - 5y^{2} + 2xy$, $\Ric^{\circ}(k,k) = 4$ and an isotropic element with $\theta_{l} = 4$ gives $r = 16$. For the vacuum wave $H = x^{3} - 3xy^{2}$, the null direction $k = \partial_{u} + \tfrac12\partial_{v} + \partial_{x}$, the flat element spanned by $e_{1} = \partial_{u} + \tfrac32\partial_{v} + 2\partial_{x}$ and $e_{2} = \partial_{y}$ (orthonormal, orthogonal to $k$ and to $l = \tfrac12\partial_{u} + \tfrac94\partial_{v} + \tfrac32\partial_{x}$, with $g(k,l) = -1$), and $k^{u} = 1$ along the element, $B^{(k)} = B^{(l)} = 0$ and $\omega = 0$ at $p$, and $R(e_{1},e_{2},e_{1},e_{2}) = R_{uyuy} = 3x$, so both rates equal $2\,k(3x) = 6$. On a radial generator of the Schwarzschild-de Sitter metric, a flat element tangent to the sphere has $\theta_{l} = \sigma_{l} = 0$, $\omega = 0$, $\Ric^{\circ}(k,k) = 0$ and $\widehat{C}_{k} = 0$ (a radial null direction is a repeated principal null direction) and $Z = 0$ by spherical symmetry, so (i) reduces to the gradient term; its value there, a change of the Euler entropy with no heat and no area change, is computed in the result filed alongside, synthesis.series-wald-entropy-changes-without-flux-on-local-horizons, and is not recomputed here.

\emph{Step 8 (consequences).} By Step 6, with $\alpha \neq 0$ the Clausius relation holds on every admissible patch through every point of a connected region, for the Jacobson-Myers form, only if the region has constant curvature; if Einstein's equations hold there, $8\pi G\,T(k,k) = \Ric^{\circ}(k,k) = 0$ and $C = 0$: no null energy and no Weyl curvature. The examples of Step 7 show the projected form fails in vacuum too. So the entropy that (H1) can use on the Einstein route does not contain the Euler term in either form. If (H1) were imposed only on flat elements ($B^{(l)} = 0$ and $\omega = 0$ at $p$), Step 6 would not apply and the constant-curvature conclusion would not follow, but $\beta = 0$ would: the gradient term alone is nonzero on the vacuum plane wave of Step 7 and on radial Schwarzschild-de Sitter generators, so the Euler coefficient must vanish already there. For the family $\frac{1}{4G}\int(\eta' + \beta\,{}^{(2)}\!R)\,dA$ (or with $P$), the Euler part contributes $\frac{\beta}{4G}r$ at first order and the area part nothing, so (H1) on a region that is not of constant curvature forces $\beta = 0$; with clause (iv) of Theorem \claimref{synthesis.c-loc-entropy-clause-curvature-independent-density}, $\eta'$ is constant. This is clause (Loc) of Premise \claimref{synthesis.c-loc-local-horizon-premise} for that family: the density depends on no curvature at the point. Densities built from other curvature components (for example $\Ric^{\circ}(k,k)$ or $R$, as in Ricci-polynomial extensions) are not treated, and (Loc) is not derived for them here. The argument needs the entropy to be a functional of the cut, so that its change between $\lambda_{0}$ and $0$ includes the change of the density; on the differences reading $\delta S = \eta\,\delta A$ with $\eta$ evaluated at $p$ there is no such term, and a curvature-dependent $\eta$ is consistent with (H1) (filed alongside, synthesis.loc-independent-of-clausius-and-killing-data).

\emph{Remarks.} (a) The rate is a total divergence on the cut: in two dimensions $\int\sqrt{h}\,{}^{(2)}\!R$ is topological, and its first variation under $\delta h_{ab} = 2B^{(k)}_{ab}\,\delta\lambda$ is $\sqrt{h}\,(2D^{a}D^{b}\sigma^{(k)}_{ab} - D^{2}\theta_{k})\,\delta\lambda$ (first variation of the two-dimensional Einstein-Hilbert density, cited from memory). So on closed cuts it integrates to zero, as Theorem \claimref{synthesis.euler-term-horizon-entropy-shift}(ii) states; on an open patch in the thin-pencil limit it does not, and the formula of (i) evaluates it at $p$ through the Codazzi equation, whose trace carries the Weyl component $C(l,k,k,e_{b})$ as well as $\Ric^{\circ}(k,e_{b})$. (b) The two forms of the Wald entropy differ off bifurcation surfaces by the extrinsic term of the Gauss equation, one of the ambiguities of Jacobson, Kang and Myers \cite{JacobsonKangMyers1994}; both fail. Clausius derivations of Lovelock equations with the Jacobson-Myers entropy exist under extra conditions on the patch (Guedens, Jacobson and Sarkar 2012; Parikh and Svesko 2018, both cited from memory and not used): as recalled here, the latter uses closed spherical cuts of a stretched light cone, on which the Euler rate integrates to zero by Remark (a), and the former restricts the local horizon beyond vanishing expansion and shear at the point. Either restriction removes some of the freedom in $\theta_{l}$, $\sigma_{l}$, $\omega$ used in Step 6, or closes the cut, so they do not conflict with this result, which concerns (H1) as the store states it: every admissible open element. In four dimensions the Euler term leaves the field equations Einstein's, and the obstruction found here is in the entropy, not in the equations. (c) Nothing here bears on the de Sitter or Schwarzschild-de Sitter entropies of closed horizons, which Theorem \claimref{synthesis.sds-clausius-fixes-area-constant} already ties to the area law with zero constant.

\emph{The check.} Exact truncated power series in $(s_{1}, s_{2}, \lambda)$ over $\mathbb{Q}$: the null normal of a polynomial element is solved as a series with prescribed $k^{u}$, the generators $X + \lambda k$ give the induced metric, its scalar curvature is computed from Christoffel symbols with series inverses and square roots, and the $\lambda$-coefficient of $\sqrt{h}\,{}^{(2)}\!R$ at $s = 0$ is the rate. The Riemann tensor of the polynomial metric and its first derivatives at $p$ are computed independently, and the projected rate is computed from the dragged frame $e_{a} + \lambda\,\partial_{a}k$. On three polynomial metrics with generic rational quadratic and cubic terms in every component, and elements with generic $B^{(l)}$, normalisation gradient and cubic data, both formulas of (i) hold exactly, every term is nonzero in some case, and the Bianchi form of the gradient term agrees. The plane-wave cases of Step 7, flat space (rate $0$ for a curved element), the independence of third- and fourth-order element data and of the normalisation where $Z = 0$, the rank statements of Step 6, and the order counting of Step 5 for three values of $\Lambda$ (including a negative one) are checked separately.
